\uuid{q2U2}
\exo7id{4875}
\titre{exo7 4875}
\auteur{quercia}
\organisation{exo7}
\datecreate{2010-03-17}
\isIndication{false}
\isCorrection{false}
\chapitre{Géométrie affine dans le plan et dans l'espace}
\sousChapitre{Applications affines}
\cpge{mpsi-mp2i.raisonnement-vocabulaire-ensembliste.c, pcsi.raisonnement-vocabulaire-ensembliste.d, ptsi.raisonnement-vocabulaire-ensembliste.d, mpsi-mp2i.applications-lineaires.b, pcsi.applications-lineaires.b, ptsi.applications-lineaires.b, mpsi-mp2i.sous-espaces-affines, ptsi.espaces-vectoriels.b, bcpst1.applications-lineaires-matrices, bcpst2.applications-lineaires-matrices.a}
\module{Géométrie}
\niveau{L2}
\difficulte{}

\contenu{
\texte{
Soit ${\cal P}$ un plan, et $f : {\cal P} \to {\cal P}$ une application
affine telle que $f^3 = \mathrm{id}$, avec $f \ne \mathrm{id}$.
}
\begin{enumerate}
    \item \question{Montrer que si $A \ne f(A)$, alors $A,f(A), f^2(A)$ sont non alignés.}
    \item \question{En déduire que $f$ est le produit de deux symétries.}
\end{enumerate}
}
